% Part: second-order-logic
% Chapter: sol-and-sets
% Section: power-of-continuum

\documentclass[../../../include/open-logic-section]{subfiles}

\begin{document}

\olfileid{sol}{set}{pow}

\olsection{连续统的势}

\begin{explain}
在二阶逻辑中，我们可以量化论域的子集，但不能量化论域子集的集合。要直接做到后者，我们需要\emph{三阶}逻辑。例如，如果我们想陈述康托尔定理——即不存在从一个集合的幂集到该集合本身的单射函数，我们可能会试图将其表述为“对每个集合~$X$ 和每个集合~$P$，如果 $P$ 是 $X$ 的幂集，则并非 $\cardle{P}{X}$”。而要说 $P$ 是 $X$ 的幂集，就需要将“$P$ 的元素恰好是~$X$ 的所有子集”形式化，即形如 $\lforall[Y][(P(Y)
  \liff Y \subseteq X)]$。问题在于 $P(Y)$：它不是二阶逻辑的公式，因为只有项才能作为像~$P$ 这样的一元关系变元的自变元。

然而，如果论域足够大，我们可以\emph{模拟}对集合的集合的量化。其想法是利用二元关系~$R$ 能将论域的元素与论域的元素相联系这一事实。给定这样一个~$R$，我们可以收集所有与某个~$x$ 有 $R$ 关联的元素：$\Setabs{y \in
  \Domain{M}}{R(x,y)}$ 就是由~$x$ “编码”的集合。反过来，如果 $Z \subseteq \Pow{\Domain{M}}$ 是~$\Domain{M}$ 的一个子集的集合，并且~$\Domain{M}$ 的元素至少和~$Z$ 中的集合一样多，那么也存在一个关系~$R \subseteq \Domain{M}^2$，使得每个 $Y \in Z$ 都能被某个~$x$ 通过~$R$ 编码。
\end{explain}

\begin{defn}
若 $R \subseteq \Domain{M}^2$，则\emph{$x$ $R$-编码} $\Setabs{y
  \in \Domain{M}}{R(x, y)}$。
\end{defn}

如果一个元素~$x \in \Domain{M}$ $R$-编码一个集合 $Z \subseteq \Domain{M}$，那么一个集合 $Y \subseteq \Domain{M}$ 就编码了一个集合的集合，即由~$Y$ 的元素所编码的那些集合。因此，一个集合 $Y$ 可以 $R$-编码 $\Pow{X}$。而 $Y$ $R$-编码 $\Pow{X}$ 当且仅当：对于每个 $Z \subseteq X$，存在某个 $x \in Y$ $R$-编码~$Z$，并且每个 $x \in Y$ 都 $R$-编码某个 $Z \subseteq X$。

\begin{prop}
公式
  \[
  \fn{Codes}(x, R, Z) \ident \lforall[y][(Z(y) \liff
    R(x, y))]
  \]
表达了 $s(x)$ $s(R)$-编码 $s(Z)$。公式
\begin{multline*}
  \fn{Pow}(Y, R, X) \ident \\
  \lforall[Z][(Z \subseteq X \lif \lexists[x][(Y(x) \land
    \fn{Codes}(x, R, Z))])] \land {} \\ \lforall[x][(Y(x) \lif
  \lforall[Z][(\fn{Codes}(x, R, Z) \lif Z \subseteq X)]]
\end{multline*}
表达了 $s(Y)$ $s(R)$-编码~$s(X)$ 的幂集，即 $s(Y)$ 的元素恰好 $s(R)$-编码了 $s(X)$ 的所有子集。
\end{prop}

\begin{explain}
利用这个技巧，我们可以通过量化子集的编码而不是子集本身，来表达关于幂集的陈述。例如，康托尔定理现在可以表述为：不存在从任何编码了~$X$ 幂集的关系的论域到~$X$ 本身的单射函数。
\end{explain}

\begin{prop}
语句
\begin{multline*}
  \lforall[X][\lforall[Y][\lforall[R][(\fn{Pow}(Y, R, X) \lif \\
      \lnot \lexists[u][(\lforall[x][\lforall[y][(\eq[u(x)][u(y)] \lif
            \eq[x][y])]] \land {}\\
        \lforall[x][(Y(x) \lif X(u(x)))])])]]]
\end{multline*}
是有效的。
\end{prop}

\begin{explain}
可数无穷集合的幂集是不可数的，因此它的势大于任何可数无穷集合的势（即 $\aleph_0$）。$\Pow{\Nat}$ 的势被称为“连续统的势”，因为它与实数线~$\Real$ 上的点等势。如果论域足够大，能够编码一个可数无穷集合的幂集，我们就可以通过断言一个集合与任何编码了势为~$\aleph_0$ 的集合~$X$ 之幂集的集合~$Y$ 等势，来表达该集合具有连续统的势。（如果论域不够大，即它不包含任何与~$\Real$ 等势的子集，那么也不存在能编码~$\Pow{X}$ 的关系。）
\end{explain}

\begin{prop}
若 $\cardle{\Real}{\Domain{M}}$，则公式
\begin{multline*}
\fn{Cont}(Y) \ident 
\lexists[X][\lexists[R][((\fn{Aleph}_0(X) \land
      \fn{Pow}(Y, R, X)) \land \\ 
      \lforall[x][\lforall[y][((Y(x) \land {} 
      Y(y) \land \lforall[z][R(x,z) \liff R(y,z)]) \lif \eq[x][y])]])]]
\end{multline*}
表达了 $\cardeq{s(Y)}{\Real}$。
\end{prop}

\begin{proof}
  $\fn{Pow}(Y, R, X)$ 表达了 $s(Y)$ $s(R)$-编码~$s(X)$ 的幂集，而 $\fn{Aleph}_0(X)$ 断言 $s(X)$ 是可数无穷的。所以 $s(Y)$ 的势至少与连续统的势一样大，尽管它可能更大（如果 $s(Y)$ 中有多个元素编码了~$X$ 的同一个子集）。这一可能性被最后一个合取枝排除，该合取枝要求经由 $s(R)$ 建立的 $s(Y)$ 的元素与 $s(Z)$ 的子集之间的关联是单射的。
\end{proof}

\begin{prop}
$\cardeq{\Domain{M}}{\Real}$ 当且仅当
\begin{multline*}
  \Sat{M}{\lexists[X][\lexists[Y][\lexists[R][(\fn{Aleph}_0(X) \land \fn{Pow}(Y, R, X) \land \\
          \lexists[u][(\lforall[x][\lforall[y][(\eq[u(x)][u(y)] \lif \eq[x][y])]] \land {} \\\lforall[y][(Y(y) \lif \lexists[x][\eq[y][u(x)]])])])]]]}.
        \end{multline*}
  \end{prop}

\begin{explain}
连续统假设是指：连续统的势是第一个不可数基数，即 $\Pow{\Nat}$ 的势是~$\aleph_1$。
\end{explain}

\begin{prop}
连续统假设为真当且仅当 \[\fn{CH} \ident
\lforall[X][(\fn{Aleph}_1(X) \liff \fn{Cont}(X))]\] 是有效的。
\end{prop}

需要注意的是，并非“$\lnot \fn{CH}$ 是有效的当且仅当连续统假设为假”。在可数论域中，既没有势为~$\aleph_1$ 的子集，也没有势为连续统的子集，因此 $\fn{CH}$ 在可数论域中总是为真。然而，我们可以给出另一个语句，它有效当且仅当连续统假设为假：

\begin{prop}
连续统假设为假当且仅当 \[\fn{NCH} \ident
\lforall[X][(\fn{Cont}(X) \lif \lexists[Y][(Y \subseteq X \land \lnot
    \fn{Count}(Y) \land \lnot \cardeq{X}{Y})])]\] 是有效的。
\end{prop}

\end{document}